\section{Сравнение классов LL и LR}

Иерархию языков, распознаваемых различными классами алгоритмов, можно представить как изображено на рисунке~\ref{fig:ll_lr_comparison}.

\begin{marginfigure}
\begin{center}
  \input{part_02_Foundations/chapter_07_ClassicalParsing/figures/06_LLvsLR/LL_LR}
\end{center}
\caption{Соотношение между классами $LL(k)$ и $LR(k)$}
\label{fig:ll_lr_comparison}
\end{marginfigure}

Из диаграммы видно, что класс языков, распознаваемых LL(k) алгоритмом уже, чем класс языков, распознаваемый LR(k) алгоритмом, при любом конечном $k$. Приведём несколько примеров.
\begin{enumerate}
\item $L = \{a^mb^nc \mid m \geq n \geq 0\} $ является LR(0), но для него не существует LL(1) грамматики.
\item $L = \{ a^n b^n + a^n c^n \mid n > 0\}$ является LR, но не LL.
\item Больше примеров можно найти в работе Джона Битти~\cite{BEATTY1980193}.
\end{enumerate}
